% PPI: Vorlage für ein Gedächtnisprotokoll
% Auf Grundlage der Vorlage von Jules Kreuer / not_a_feature (2021).
% Mit pdfLaTeX kompilieren und die fertige PDF-Datei hochladen.
\documentclass[11pt,a4paper]{article}
\usepackage[T1]{fontenc}
\usepackage[utf8]{inputenc}
\usepackage[ngerman]{babel}
\usepackage[margin=2.5cm]{geometry}
\usepackage{amsmath,amssymb}
\setlength{\parindent}{0pt}
\setlength{\parskip}{0.6em}

\begin{document}
\begin{center}
  {\Large\bfseries Gedächtnisprotokoll}\\[0.5em]
  {\large [Name der Vorlesung]}
\end{center}

\begin{description}
  \item[Prüfende Person:] [Name]
  \item[Semester:] [z.\,B. Winter 2026/27]
  \item[Prüfungsdatum:] [TT.MM.JJJJ]
  \item[Prüfungsart:] [Mündlich / schriftlich; Haupt- / Nachprüfung]
  \item[Dauer:] [Minuten]
  \item[Hilfsmittel:] [Erlaubte Hilfsmittel]
\end{description}

\section*{Frage / Aufgabe 1}
[Frage oder Aufgabenstellung aus dem Gedächtnis beschreiben.]

\subsection*{Antwort / Lösungsweg}
[Eigene Antwort oder den Lösungsweg ergänzen. Unsichere Angaben kennzeichnen.]

\section*{Frage / Aufgabe 2}
[Weitere Frage oder Aufgabe beschreiben.]

\subsection*{Antwort / Lösungsweg}
[Antwort oder Lösungsweg ergänzen.]

\section*{Anmerkungen}
[Optional: Ablauf, Schwerpunkte oder Hinweise zur Vorbereitung.]
\end{document}
